\section{Distributed Execution: Proofs and Details}
\label{app:proofs-sec4}

This appendix proves the results of Section~\ref{sec:distributed-execution}, supplies the
\commitment bases behind Section~\ref{sec:systems}, works out the bases of two-phase
locking and multiversion timestamp ordering, and extends Complete CAP to deployments and
to explicit relations that the system chooses.

\subsection{Proofs}

\begin{proof}[Proof of Proposition~\ref{prop:distributed-consistency}]
A cut is confluent at $j$ exactly when its participant runs are
(\S\ref{sec:distributed-consistency}).  First assume runs are consistent.  The
participant runs at any cut form a finite family of runs, so histories are
consistent.  Now assume histories are consistent.  Every finite family of runs
is the family of participant runs, at the final cut, of a message-free history
with one participant per run, so runs are consistent.
\end{proof}

\paragraph{Exchanging outputs.}
\label{app:exchanging-outputs}
Suppose participants share their \emph{outputs} in the data $d$ of messages
(\S\ref{sec:distributed-consistency}).  A participant that receives an
output $o$ its sender has emitted can do either or both of two things.  If
$o\in R(i)$ at its own current input $i$, it can emit $o$ itself, an ordinary
output event of Definition~\ref{def:distributed-history}.  It can also
silently honor $o$ by restricting its table to $R[\mathrm{out}\OutputGe o]$.
The restricted table may have no row at the participant's current input, so it
is not always a specification state.  Silent honoring is therefore not a
specification event, and we reason about it on tables.

Silent honoring never changes whether a cut is confluent at a common \indesc
$j$.  It cannot make a cut confluent at $j$, because it only removes rows, and
confluence asks for a row that survives everywhere
(Proposition~\ref{prop:confluence}).  It cannot make a cut non-confluent at
$j$ either.  A row that survives at $j$ already
\outrel every output emitted in the cut (Lemma~\ref{lem:run-residual}),
including $o$, so the restriction keeps the row.  The confluence of a cut
therefore depends on which outputs were emitted, not on who has heard about
them.

\begin{proof}[Proof of Proposition~\ref{prop:system-confluence}]
(1) By Lemma~\ref{lem:canonical-equivalent}, $R_r\subseteq\widehat R_r$.
Rows are never added, so a row that a \commitment event of $\rho_r$ removed is
not in $R_r$; hence $R_r\subseteq\widehat R_r\setminus\Delta(\rho_r)$.
Conversely, let $(k,u)\in\widehat R_r\setminus\Delta(\rho_r)$.  Both runs
start at the same canonical state, which contains $(k,u)$.  Consider an event of
$\rho_r$ at a state containing $(k,u)$.  If it is $\mathsf{input}(j)$, the
canonical run also takes $\mathsf{input}(j)$ and keeps $(k,u)$, so
$j\InputOrd k$ and the event keeps $(k,u)$.  If it is $\mathsf{output}(o)$,
then likewise $o\OutputOrd u$ and the event keeps $(k,u)$.  If it is a
\commitment event, it keeps $(k,u)$ because $(k,u)\notin\Delta(\rho_r)$.  By
induction on the events, $(k,u)\in R_r$.

(2) By~(1), $(k,u)\in R_r$ exactly when $(k,u)\in\widehat R_r$ and
$(k,u)\notin\Delta(\rho_r)$, so it suffices to show that the runs are confluent
at $k$ exactly when some row $(k,u)$ lies in every $R_r$.  First assume
continuations bring every run to a common state $(k,R)$.  That state has a row
$(k,u)$, and events never add rows, so $(k,u)\in R_r$ for every $r$.  Now
assume $(k,u)\in R_r$ for every $r$.  Since $i_r\InputOrd k$, the event
$\mathsf{input}(k)$ is legal at $(i_r,R_r)$ and keeps $(k,u)$, so
$\mathsf{output}(u)$ is then enabled.  It ends at $(k,R'_r)$ with
$R'_r=R_r[\mathrm{in}\InputGe k,\ \mathrm{out}\OutputGe u]$.  If $R'_r$ has no
row other than $(k,u)$, the run has reached $(k,\{(k,u)\})$.  Otherwise
consider the partial function that maps each state with current input $k$
whose table contains $(k,u)$ and some other row, and whose rows at $k$ all
\outrelbase $u$, to $(k,\{(k,u)\})$.  It keeps the current input and removes a
nonempty set of rows.  A removed row at $k$ is some $(k,x)$ with
$u\OutputOrd x$ and $x\neq u$, so $x\not\OutputOrd u$; a removed row at
another input has no remaining row at its input.  So the function is a
\commitment (Definition~\ref{def:commitment}), and a \commitment event that
applies it ends at $(k,\{(k,u)\})$.  Every run therefore reaches
$(k,\{(k,u)\})$.
\end{proof}

\begin{proof}[Proof of Proposition~\ref{prop:commitment-conflict}]
Events only remove rows, so $(k,x)$ is present at $\sigma_P(i)$.  Thus $k$
\inrel $i$ and $x\in\fiber{k}$, so $\mathsf{input}(k)$ and then
$\mathsf{output}(x)$ are enabled from $\sigma_P(i)$, and a canonical second
run exists.  Under any basis, after $\mathsf{output}(x)$ every row at $k$
\outrel $x$.  The \commitment leaves no such row at $k$
(Equation~\ref{eq:commitment-unsubsumed}), and if it removed every row at $k$,
there is none at all.  Later events only remove rows, so no row at $k$
survives in both runs, and by Proposition~\ref{prop:system-confluence} the
runs are not confluent at $k$.
\end{proof}

\begin{proof}[Proof of Theorem~\ref{thm:complete-cap}]
The canonical system for $E$ has the canonical basis and no \commitment
events.  Its strategies emit any output $o$ of $E$ in the current fiber by
the canonical output event $\mathsf{output}(o)$ of
Section~\ref{sec:complete-calm}, which applies the \commitment $x_o$ from the
canonical basis (Definition~\ref{def:run-under-basis}).  So the system has
Availability for $E$ whatever the spec.

First assume $\Spec$ satisfies Lifted Monotonicity for $E$.  Every output the
canonical system emits is the output of a row of $E$ whose input lies below
the current one.  So at a common \indesc $j$ of the current inputs, the
emitted outputs have a common \outnoun in $\fiber{j}$.  The system emits only
canonical output events, so Lemma~\ref{lem:reconciliation} makes this
confluence at $j$, and Partition Consistency follows.

Now assume some system, such as the canonical system for $E$, has
Availability for $E$ and Partition Consistency.  Let
$(i_1,o_1),\ldots,(i_n,o_n)$ be rows of $E$ with a common \indesc $j$.  We
find some $u\in\fiber{j}$ that \outrel every $o_r$.  If $n=0$, then $j$
appears in some row $(j,u)$, since $\Spec$ is total.  Otherwise, take one
participant $p_r$ per row, starting at input $i_r$, each on its own side and
running without messages.  A partitioned history has at least two sides, so if
$n=1$ we add a second participant for that row on another side.  By
Availability for $E$, each $p_r$ can emit $o_r$.  By Partition Consistency,
the participants then reach one state $(j,R)$, which has some row $(j,u)$
(Definition~\ref{def:residual-state}).  Rows are never added, so $(j,u)$ was
already in each $p_r$'s table just after it emitted $o_r$.  Emitting an output
keeps only rows whose output \outrel it
(Definition~\ref{def:lifted-output-event}), so $u$ \outrel every $o_r$, and
Lifted Monotonicity for $E$ follows.
\end{proof}

\begin{proof}[Proof of Corollary~\ref{cor:complete-cap-any-system}]
The converse argument in the proof of Theorem~\ref{thm:complete-cap} applies
to any system with Availability for $E$ and Partition Consistency, so $\Spec$
satisfies Lifted Monotonicity for $E$.  By Theorem~\ref{thm:complete-cap},
$\Spec$ escapes CAP for $E$.
\end{proof}

\subsection{\Commitment bases}
\label{app:bases}\label{app:commitment-bases}

This subsection shows that the canonical \commitments of
Sections~\ref{sec:functional-calm}--\ref{sec:complete-calm} are \commitments
(Definition~\ref{def:commitment}) and form a canonical basis, and proves
Observable equivalence (Lemma~\ref{lem:canonical-equivalent}).  Because
\commitments are partial functions, a basis may offer a \commitment only at the
states where its own precondition holds.

\begin{definition}[Canonical basis]
\label{def:canonical-commitment-basis}
The \emph{canonical \commitment basis} collects the canonical \commitments of
Definition~\ref{def:canonical-commitment},
\begin{equation}
  \mathcal B_P
    =\{x_o:o\in O,\ \operatorname{dom}(x_o)\neq\varnothing\},
  \label{eq:canonical-commitment-basis-relational}
\end{equation}
where $x_o$ is defined at the states with $o\in R(i)$ and $\Excluded_o\neq\varnothing$.
\end{definition}

\paragraph{The canonical \commitment is a \commitment.}
Definition~\ref{def:commitment} asks that $x_o$ keep the current input, which
it does, remove a nonempty set of rows, which $\Excluded_o\neq\varnothing$
ensures, and remove no row that a kept row at the same input \outrel
(Equation~\ref{eq:commitment-unsubsumed}).  Only the last needs an argument.
Suppose, for contradiction, that a removed row $(k,x)\in\Excluded_o$ is
\outrelpast by a row $(k,y)$ that $x_o$ leaves.  Either $y$ \outrel $o$, or
$y\in\Subsumed_o$ and is \outrelpast by a row \outrelpart $o$.  By
transitivity, $x$ is then \outrelpast by a row \outrelpart $o$, so
$x\in\Subsumed_o$, a contradiction.

Canonical runs (Definition~\ref{def:run}) are exactly the runs under the
canonical basis $\mathcal B_P$ of Definition~\ref{def:canonical-commitment-basis}
with no \commitment events.

\begin{proof}[Proof of Lemma~\ref{lem:canonical-equivalent}]
We drop the \commitment events from a run under $\mathcal B$.  By induction on
the events, the result is a canonical run whose table $\widehat R$ contains
the original's table $R$ after each event; both start at the same canonical
state.  A \commitment event keeps the current input and only removes rows, so
dropping it preserves $R\subseteq\widehat R$.  An input event
$\mathsf{input}(j)$ legal at $R$ is legal at $\widehat R$, and keeps
$R[\mathrm{in}\InputGe j]\subseteq\widehat R[\mathrm{in}\InputGe j]$.  An
output event $\mathsf{output}(o)$ at $R$ can be issued at $\widehat R$, since
$o\in R(i)\subseteq\widehat R(i)$, with $x_o$ from the canonical basis if it
applies one, and keeps
$R[\mathrm{out}\OutputGe o]\subseteq\widehat R[\mathrm{out}\OutputGe o]$.
\end{proof}

\begin{corollary}[Reconciliation under any basis]
\label{cor:reconciliation-any-basis}
Let runs under any \commitment bases end at $(i_r,R_r)$, and let $j$ be a
common \indesc of the $i_r$.  If some continuations, whose output and \commitment events
may apply any \commitments, bring them to a common state with current input $j$,
then some $u\in\fiber{j}$ \outrel every output emitted on every run.
\end{corollary}

\begin{proof}
The common state has a row $(j,u)$, and events never add rows, so
$(j,u)\in R_r$.  By Lemma~\ref{lem:canonical-equivalent}, $R_r$ is contained in
the table of a canonical run with the same inputs and outputs, and
Lemma~\ref{lem:run-residual} makes $u$ \outrel the outputs of run $r$.
\end{proof}

The converse of Lemma~\ref{lem:canonical-equivalent} fails: a canonical run
need not be observably equivalent to any run under a particular system's
basis.  That basis may lack the \commitment that an output event requires,
which disables that output.  Looking only at the runs that one system produces
could therefore hide non-monotonicity, and hide a canonical run that conflicts
with a \commitment.  We instead judge consistency, and the effect of
\commitments, against the canonical runs of the declared spec.  A behavioral
restriction narrows that comparison only when it is part of the declared
contract.

\subsection{\Commitment bases for two-phase locking and multiversion timestamp ordering}
\label{app:cc-bases}

This subsection gives in full the serializability spec of
Sections~\ref{sec:ambiguity-story} and~\ref{sec:systems}, and the \commitment
bases of two-phase locking (2PL) and multiversion timestamp ordering (MVTO).
Schedulers report results while transactions are still running, such as a
read result before its reader commits, so an output may leave transactions
undecided.

\paragraph{Specification.}
An input is a finite set of transactions, each given by the sequence of
operations that its client has issued so far: reads and writes of keys,
possibly followed by a request to commit.  One input precedes another when
every transaction of the first is in the second and its sequence in the first
is a prefix of its sequence in the second; once a transaction has requested
commit, its sequence no longer grows.  An input does not order the operations
of different transactions.  Their arrival order is the order of input events
in a run, which the environment controls.

An output has three parts: the committed transactions, in a serial order; the
aborted transactions; and the read results reported so far.  A read result
says which version the $k$-th operation of a transaction, a read of a key $x$,
returned: the write of $x$ by some transaction (its last write of $x$), or the
initial version.  A transaction that is
neither committed nor aborted is \emph{undecided}, and its reads may already be
reported.  One output \outrel another when it extends all three parts.

An output is \emph{complete} at an input when it commits or aborts every
transaction of the input, and not both, and commits only transactions whose
commit request is in the input.  A complete output is permitted when it meets two
conditions.  First, each reported read result is for a read in the input and
returns the initial version or the write of its key by a transaction of the
input.  Second, each read of a key $x$ by a committed transaction $T$ returns
$T$'s own write of $x$ if $T$ writes $x$ earlier in its sequence.  Otherwise it
returns the write of $x$ by the last transaction before $T$ in the serial order
that writes $x$, or the initial version if none does.  Reads by aborted
transactions are otherwise unconstrained.  Any output is permitted when some permitted complete output
\outrel it, that is, when it can still be completed.  The figures of
Sections~\ref{sec:ambiguity-story}--\ref{sec:systems} show only complete
outputs without aborts or read results, as serial orders, at inputs in which
every transaction has requested commit.

We write an output as $(L,X,\rho)$, where $L$ is the serial order, a total
order on the set $\mathrm{dom}\,L$ of committed transactions, $X$ is the set of
aborted ones, and $\rho$ is the set of reported read results.  Then
\[
  \begin{aligned}
  (L,X,\rho)\OutputOrd(L',X',\rho')
  \iff{}& X\subseteq X',\ \ \mathrm{dom}\,L\subseteq\mathrm{dom}\,L',\\
         & L'|_{\mathrm{dom}\,L}=L,\ \ \rho\subseteq\rho'.
  \end{aligned}
\]
Every input has a permitted row: the one that aborts every transaction and
reports nothing.  Let $\{A,B\}$ be an input in which $A$ and $B$ have each
requested commit.  The rows $(\{A,B\},([A,B],\varnothing,\varnothing))$ and
$(\{A,B\},([B,A],\varnothing,\varnothing))$ of
Figure~\ref{fig:ambiguity-frames}(f) have no common \outnoun, since it would
order $A$ both before and after $B$.  So the spec is not monotone.  The
\emph{abort-free} spec permits a row when it can be completed with
$X=\varnothing$.  For this spec, a complete output instead commits exactly the
transactions whose commit request is in the input and leaves the others
undecided.  The same two rows show that it is not monotone either.

\paragraph{\Commitments.}
Each \commitment is a condition on complete outputs.  Applied at a state
$(i,R)$, it keeps the rows of $R$ that can still be completed, within $R$, by a
complete row satisfying the condition, and removes the rest.  It is defined
when it keeps some row at $i$ and removes some row.
\begin{itemize}
  \item $\mathsf{order}(A,B)$, for distinct $A,B\in i$: $B$ is not committed
        before $A$, that is, not $B<_L A$.
  \item $\mathsf{TxCommit}(A)$, for $A$ whose commit request is in $i$:
        $A\notin X$.\quad
        $\mathsf{abort}(A)$, for $A\in i$: $A\notin\mathrm{dom}\,L$.
  \item $\mathsf{stamp}(A)$, for $A\in i\setminus Z$, given the set $Z\subseteq
        i$ of transactions stamped before $A$: no committed transaction outside
        $Z\cup\{A\}$ precedes $A$ in $L$, and no committed transaction of $Z$
        follows $A$.
  \item $\mathsf{read}(A,k,v)$, for $A\in i$ whose $k$-th operation in $i$ is a
        read: the output reports that this read returned the version $v$.
\end{itemize}
Each is a \commitment (Definition~\ref{def:commitment}).  In particular, no
removed row is \outrelpast by a kept one, as
Equation~\ref{eq:commitment-unsubsumed} requires: if a kept row \outrel another
row at the same input, the complete row that keeps the first also \outrel the
second, so the second is kept too.  The first four conditions concern only $L$
and $X$, and for them a complete row is kept exactly when it satisfies the
condition, since two permitted complete outputs at one input, one
\outrelpart the other, have the same serial order and abort set.

The bases are
$\mathcal B_{\mathrm{2PL}}=\{\mathsf{order},\mathsf{TxCommit},\mathsf{abort},\mathsf{read}\}$
and
$\mathcal B_{\mathrm{MVTO}}=\{\mathsf{stamp},\mathsf{TxCommit},\mathsf{abort},\mathsf{read}\}$,
over all transactions, positions, and versions.  They share
$\mathsf{TxCommit}$, $\mathsf{abort}$, and $\mathsf{read}$, and differ in how
they order.  Let $\{A,B,C\}$ be an input in which $A$, $B$, and $C$ have each
requested commit, after any operations.  At its canonical state, no \commitment
of $\mathcal B_{\mathrm{MVTO}}$ has the effect of $\mathsf{order}(A,B)$, and no
\commitment of $\mathcal B_{\mathrm{2PL}}$ has the effect of $\mathsf{stamp}(B)$
with $Z=\varnothing$.

A scheduler applies its \commitments in \commitment events
$\mathsf{apply}(\gamma)$ (Definition~\ref{def:commitment-event}) and reports
results by output events $\mathsf{output}(o)$
(Definition~\ref{def:run-under-basis}).  Each protocol step below responds to
the arrival of an operation, an input event, and says which \commitments it
applies.  Every state that a run under either basis reaches has two
properties: each row has a complete \outnoun in the table at its input, and
any output between two rows of the table at one input is itself a row of the
table.  So a scheduler can report a read by first applying
$\mathsf{read}(A,k,v)$ in a \commitment event, after which emitting the result
is a legal output.

\paragraph{2PL.}
When an operation of $B$ arrives, 2PL requests a lock on its key, shared for a
read and exclusive for a write.  If $A$ holds that lock in a conflicting mode,
$B$ waits, and the scheduler applies $\mathsf{order}(A,B)$ in a \commitment
event and reports nothing.  If every remaining complete row that commits both
already places $A$ before $B$, then $\mathsf{order}(A,B)$ would remove
nothing, so it is undefined and the scheduler applies nothing.  This happens
when an earlier wait has applied $\mathsf{order}(A,B)$, or when a reported
read result fixes the order, as when $A$ has read the initial version of the
key that $B$ writes.  On complete rows without aborts, $\mathsf{order}(A,B)$
keeps exactly those with $A$ before $B$ (Figure~\ref{fig:cc-commitments}(h)).
At the canonical state of $\{A,B,C\}$ it removes $[B,A,C]$ and keeps
$[A,B,C]$ and $[C,A,B]$.  So it is a \commitment, and it is not the canonical
\commitment $x_o$ of any output $o$.  When a read is granted its lock, the
scheduler applies $\mathsf{read}$ and reports the result.  When a commit
request arrives, the scheduler applies $\mathsf{TxCommit}$ and releases the
transaction's locks.

A deadlock between $A$ and $B$ arises when $B$ waits for a lock that $A$ holds
and an operation of $A$ then waits for a lock that $B$ holds.  The scheduler
has then applied $\mathsf{order}(A,B)$ and later $\mathsf{order}(B,A)$, so no
remaining row, at any input, commits both $A$ and $B$.  Every remaining
complete row aborts one of them, and resolving the deadlock means aborting
one.  Once the scheduler applies $\mathsf{TxCommit}(A)$, every remaining
complete row aborts $B$, and $\mathsf{TxCommit}(B)$ is never again defined.
Under the abort-free spec, once both orderings have been applied, no remaining
row is at an input in which both $A$ and $B$ have requested commit.  The
system never accepts both requests, so at least one of $A$ and $B$ waits
forever.  Whether a deadlock can be resolved thus depends on which of the two
specs the system implements.

\paragraph{MVTO.}
We follow Bernstein and Goodman~\cite{bernstein1981concurrency}.  When the
first operation of a transaction $A$ arrives, $A$ receives a timestamp above
all issued ones, and the scheduler applies $\mathsf{stamp}(A)$, with $Z$ the
transactions stamped before $A$, in a \commitment event.  On complete rows
without aborts, $\mathsf{stamp}(B)$ with $Z=\varnothing$ keeps exactly those
with $B$ first (Figure~\ref{fig:cc-commitments}(i)).  At the canonical state
of $\{A,B,C\}$ it is a \commitment, and it is not the canonical \commitment
$x_o$ of any output.  A read of $x$ by $A$ returns $A$'s own write of $x$ if
$A$ has written $x$, and otherwise the committed version of $x$ with the
largest timestamp below $A$'s.  The scheduler applies $\mathsf{read}$ and
reports the result.  A write of $x$ by $A$ is rejected, and $A$ aborted, if a
transaction stamped after $A$ has already read a version of $x$ stamped before
$A$.  The scheduler then applies $\mathsf{abort}(A)$ and reports the abort.
When the commit request of $A$ arrives, the scheduler applies
$\mathsf{TxCommit}(A)$.

\emph{What reporting a read excludes.}  Suppose $B$ is stamped after $A$, and
$B$'s read of $x$ returned the initial version or the write of $x$ by a
transaction stamped before $A$.  Consider an input in which $A$ writes $x$.
There, a complete row that satisfies both stamps and commits both $A$ and $B$
places $A$ between that version and $B$.  In that row the read must return
$A$'s write or a later one, so the row does not permit the reported read.
Hence, after the two stamps, no row at an input in which $A$ writes $x$ both
reports the read and commits $A$ and $B$.  Once $\mathsf{read}$ has been
applied, no remaining row at such an input commits both, whether it reports
the read or not.  Reporting the read therefore rules out one possibility: that
$A$ writes $x$ and both $A$ and $B$ commit.  MVTO reports the read when it
arrives, whether or not $A$'s write has arrived.

The condition on the version, that $B$ read a version older than $A$, is
needed.  Consider an input of three transactions $A$, $W$, $B$, stamped in
that order, each having requested commit, with $A$ and $W$ writing $x$ and $B$
only reading it.  The complete row that commits all three in stamp order and
reports that $B$ read $W$'s write is permitted and satisfies every stamp.
There $B$ read a version written after $A$, and all three commit.

\emph{When the abort is entailed.}  Which of $A$ and $B$ aborts depends on
whether $B$ is decided.  Suppose $A$'s write of $x$ has arrived.  If every
remaining complete row commits $B$, for instance after $\mathsf{TxCommit}(B)$,
then every remaining complete row aborts $A$.  An output that reports the read
and that every remaining row \outrel, with $A$'s abort added, is then a legal
output.  So the spec entails the abort.  While $B$ is undecided, the spec does
not entail it.

Consider a run in which $A$'s first operation, a read of $x$, arrives, then
$B$'s read of $x$, and then $A$'s write of $x$, with further transactions free
to arrive.  The scheduler applies $\mathsf{stamp}(A)$ and $\mathsf{stamp}(B)$
in \commitment events and reports that $B$ read the initial version of $x$.
Reporting the read is not a legal output on its own, since at an input in
which $A$ writes $x$, the row that commits $A$ and then $B$ cannot be
completed with it.  Once the scheduler applies $\mathsf{read}$ in a
\commitment event, reporting the read is a legal output.  After $A$'s write
arrives, two completions remain at the input in which both then request
commit: one commits $B$ and aborts $A$, and the other commits $A$ and aborts
$B$.  Emitting $A$'s abort is then not a legal output, and
$\mathsf{abort}(A)$ and $\mathsf{abort}(B)$ are both defined and keep
different completions.  MVTO's write rule aborts $A$ at once: the scheduler
applies $\mathsf{abort}(A)$ in a \commitment event, after which reporting the
abort is a legal output.  That abort is the scheduler's choice between two
outcomes the spec still permits.  If the scheduler instead waits until $B$'s
commit request arrives and applies $\mathsf{TxCommit}(B)$, $A$'s abort becomes
a legal output.

\paragraph{Comparison.}
As operations arrive, 2PL applies \commitments that fix the relative order of
conflicting transactions only.  MVTO, when a transaction's first operation
arrives, applies a \commitment that fixes its position against all present and
future transactions.  Both report reads through $\mathsf{read}$.  Under MVTO,
reporting a read can itself rule out that two transactions both commit, and
the write rule settles this by an abort that the spec forces only after the
reader commits.

\subsection{Deployments: one explicit relation per participant}
\label{app:deployment}

Theorem~\ref{thm:complete-cap} fixes one explicit relation $E$ for all
participants.  A deployment fixes one for each participant: the rows that
participant may emit.  It models a system that assigns authority, such as a
key owned by one participant or a replica that answers only its own reads.
Assigning rows to participants removes no row from the spec, and
\commitments are still compared with all of the spec's canonical runs.

\begin{definition}[Deployment, emitted rows, partition]
\label{def:deployment}\label{def:emitted-rows}
A \emph{deployment} over a finite participant set $\mathcal N$ assigns to each $p\in\mathcal N$
an explicit relation $E_p\subseteq P$, the rows $p$ may emit: $p$ performs an output event
$\mathsf{output}(o)$ (Definition~\ref{def:run-under-basis}) at current input
$i$ only if $(i,o)\in E_p$.  The \emph{emitted rows} of a run are the
pairs $(i,o)$ such that the run emits $o$ at current input $i$; the emitted
rows of a system are those of its participants' runs in the histories it
produces.  A \emph{partition} $\Pi$ of $\mathcal N$ divides $\mathcal N$ into
two or more nonempty sides, as in a partitioned history
(Definition~\ref{def:partitioned-history}).
\end{definition}

\begin{definition}[Participant run family, silent participant, separable run family]
\label{def:participant-run-family}
A \emph{participant run family} assigns to each of finitely many participants
$p\in\mathcal N$ an initial input and a canonical run whose emitted rows lie in $E_p$.
A participant whose run emits no output is \emph{silent}.  The family is
\emph{separable under $\Pi$} when either participants on at least two
sides emit, or exactly one distinct row is emitted and some silent
participant lies on another side.
\end{definition}

\begin{definition}[Monotone across partitions]
\label{def:monotone-across-partitions}
Relative to a deployment, $\Spec$ is \emph{monotone across $\Pi$} when, for
every participant run family separable under $\Pi$ and every common \indesc $j$
of its participants' current inputs, some $u\in\fiber{j}$ \outrel every
emitted output.  It is \emph{monotone across partitions} when it is monotone
across every partition of $\mathcal N$.
\end{definition}

Monotonicity across partitions plays the role of Lifted Monotonicity for $E$
(Definition~\ref{def:lm-for-A}), restricted to the families that a partition
can separate.  If $\Spec$ satisfies Lifted Monotonicity for a relation that
contains every $E_p$, it is monotone across partitions, since the emitted rows
of a family are rows of that relation at inputs below $j$.  So Lifted
Monotonicity for $E$ suffices for the uniform deployment $E_p=E$, and Lifted
Monotonicity suffices for every deployment.

The converse can fail, because a deployment can keep conflicting rows at one
participant.  Rows authorized only to $p$ can be emitted only by $p$'s run,
and the outputs that run emits form a chain, since each output event keeps
only rows whose outputs extend it.  So a family whose emitted rows are all
authorized only to $p$ is separable only when it emits a single row.  For
example, let $\mathcal N=\{p,q\}$, let every fiber be $\{a,b\}$ ordered by
equality, and let $E_p=P$ and $E_q=\varnothing$, so that $p$ alone names a
leader.  A run of $p$ emits at most one distinct output, and $q$ never emits.
Every separable family therefore consists of one row $(i,a)$ or $(i,b)$ beside
a silent $q$.  That output lies in $\fiber{j}$ for every \indesc $j$ of $i$.
So $\Spec$ is monotone across partitions.  Yet $(i,a)$ and $(i,b)$ have no
common \outnoun, and Lifted Monotonicity fails.

The same happens whenever two conditions hold.  First, each authorized row has
an \outnoun at every \indesc of its input.  Second, every set of authorized
rows that lacks a common \outnoun at some common \indesc is authorized only to
a single participant.  The causal register of Appendix~\ref{app:registers} is
monotone across partitions without Lifted Monotonicity and has no single
authority.  It need not satisfy these conditions, since a set of rows lacking
a common \outnoun may also contain rows of other participants.

The converse does hold when every participant may emit every row.  Suppose
$|\mathcal N|\ge2$ and $E_p=P$ for every $p$.  Then every pair of rows can be
split across a partition, and every single row can be paired with a silent
participant at the common \indesc.  Pairs of rows suffice for the
common-extension condition (Appendix~\ref{app:pairs-suffice}), so
monotonicity across partitions is then Lifted Monotonicity.

\begin{definition}[Deployment Availability]
\label{def:deployment-availability}
Relative to a deployment, a system has \emph{Deployment Availability} when, in
every partitioned history over $\mathcal N$ it produces, each participant $p$
can emit every output in $R(i)\cap E_p(i)$, where $(i,R)$ is $p$'s final state
(Definition~\ref{def:local-continuation}).
\end{definition}

This is Availability for $E_p$, one relation per participant.  For the uniform
deployment $E_p=E$ it is Availability for $E$ on histories over $\mathcal N$.

A deployment is \emph{single-authority} when $E_q=\varnothing$ for every
participant $q$ but one, as in the leader example above.  Every separable
family is then one row emitted by the single authority beside a silent
participant, and monotonicity across partitions says only that each such row
has an \outnoun at every \indesc of its input.  Placement-sensitive models
such as relational-transducer CALM~\cite{ameloot2013relational,ameloot2015weaker}
also cover this case.

A deployment keeps the conflicting choice from being made at another
participant, but it does not remove the semantic conflict: a canonical run
still conflicts with the authority's \commitment
(Proposition~\ref{prop:commitment-conflict}).  For example, another canonical
run of the spec conflicts with a lock manager's choice even when only the
manager is authorized to emit the result.  A deployment also does not restrict
which inputs a participant receives: the environment may deliver any input to
any participant.

\begin{proposition}[Partition compatibility under a deployment]
\label{thm:complete-cap-deployment}
Fix a deployment.  The following are equivalent.
\begin{enumerate}
  \item For every partition $\Pi$, every participant run family separable
        under $\Pi$ is confluent at every common \indesc $j$ of its
        participants' current inputs.
  \item $\Spec$ is monotone across partitions.
\end{enumerate}
\end{proposition}

\begin{proof}
Fix $\Pi$, a separable family, and $j$, and let $X$ be the set of emitted
outputs.  By Lemma~\ref{lem:reconciliation}, the family is confluent at $j$ if
and only if some $u\in\fiber{j}$ \outrel every output in $X$.  The two
conditions thus agree family by family.
\end{proof}

The proposition concerns canonical runs and does not use Deployment
Availability.  It also applies to the runs of a system, since by
Lemma~\ref{lem:canonical-equivalent} each such run is observably equivalent to
a canonical run with the same emitted rows, which therefore lie in $E_p$.  So
if $\Spec$ is monotone across partitions, a separable family of system runs in
a partitioned history, emitting as Deployment Availability permits, has
emitted outputs with a common \outnoun $u\in\fiber{j}$ at every common
\indesc $j$.  The deployment need not authorize any participant to emit $u$.
Confluence through authorized output events would also require that
authority, a condition we do not impose.

With $|\mathcal N|\ge2$ and $E_p=P$ for every $p$, item~(2) is Lifted
Monotonicity.  The proposition is then the step behind
Theorem~\ref{thm:complete-cap} with $E=P$: separable families are confluent at
every common \indesc exactly when Lifted Monotonicity holds.

For systems, a deployment has its own form of Complete CAP, which rests on one
more fact about families: with two or more participants, monotonicity across
partitions already covers every family, separable or not.

\begin{lemma}[Compatibility of participant run families]
\label{lem:deployment-compatible}
Fix a deployment with $|\mathcal N|\ge2$.  Then $\Spec$ is monotone across
partitions if and only if, for every participant run family and every common
\indesc $j$ of its participants' current inputs, some $u\in\fiber{j}$
\outrel every emitted output.
\end{lemma}

\begin{proof}[Proof sketch]
The condition covers every family, separable or not, so it implies
monotonicity across partitions.  For the other direction, assume $\Spec$ is
monotone across partitions.  If participants $p$ and $q$ both emit, a
partition with $p$ on one side and $q$ on the other makes the family
separable, and the separable case applies.  If no participant emits, any
$u\in\fiber{j}$ will do.  Otherwise exactly one participant $p$ emits.  Its
outputs form a chain, so the last of them \outrel the rest.  Consider $p$'s
last row beside one other participant, placed on another side and kept
silent.  This is a separable family at the same $j$, and the common \outnoun
for that row \outrel every output of $p$.
\end{proof}

The \emph{canonical system} of a deployment is the system whose participants
follow canonical runs, each participant $p$ emitting any output of $E_p$ at
its current input that extends its earlier outputs.

\begin{proposition}[Complete CAP for a deployment]
\label{prop:deployment-complete-cap}
Fix a deployment with $|\mathcal N|\ge2$.  Some system has Deployment
Availability and Partition Consistency on histories over $\mathcal N$ if and
only if $\Spec$ is monotone across partitions.  When it is, the deployment's
canonical system is such a system.
\end{proposition}

\begin{proof}[Proof sketch]
First assume $\Spec$ is monotone across partitions.  The canonical system has
Deployment Availability, since each participant can emit any authorized output
that its state allows.  Its participants' runs in a partitioned history form a
participant run family, so by Lemma~\ref{lem:deployment-compatible} their
emitted outputs have a common \outnoun in $\fiber{j}$.  By
Lemma~\ref{lem:reconciliation} the runs are then confluent at $j$, so
Partition Consistency follows.

Now assume some system has both properties.  Consider a participant run
family and a common \indesc $j$.  Each participant's emitted outputs form a
chain, so it suffices to reach each participant's last emitted row.  We start
each participant alone on its own side, at the input of that row.  Deployment
Availability lets it emit the row's output.  Partition Consistency then gives
the emitted outputs a common \outnoun in $\fiber{j}$, so $\Spec$ is monotone
across partitions.
\end{proof}

\subsection{Choosing the explicit relation}
\label{app:cap-response}

Theorem~\ref{thm:complete-cap} fixes the explicit relation $E$ in advance.
Gilbert and Lynch ask only that every request receive some
response~\cite{gilbert2002cap}, and leave the system to choose which.  Here
the system chooses its explicit relation among declared acceptable responses.
The result is Theorem~\ref{thm:complete-cap} applied to the chosen relation
(Theorem~\ref{thm:complete-cap-response}), and it yields the Gilbert--Lynch
theorem with precedence taken from real time (Corollary~\ref{cor:gilbert-lynch}).
Throughout, we fix an \spectype $\Spec$ and the partitioned histories of
Section~\ref{sec:complete-cap}; the environment supplies any input to any
participant and resolves the strategies' nondeterministic choices.  Emitted
rows are as in Definition~\ref{def:emitted-rows}.

\subsubsection{Obligations, availability, and locality}

\begin{definition}[Response obligation, answering]
\label{def:response-obligation}
A \emph{response obligation} assigns to each participant $p$ a relation
$Q_p\subseteq P$ that is total on $I$ and upward closed within fibers: if
$(i,o)\in Q_p$, $(i,u)\in P$, and $o\OutputOrd u$, then $(i,u)\in Q_p$.  An
output in $Q_p(i)$ \emph{answers} at $i$ for $p$.
\end{definition}

\begin{definition}[Uniform obligation]
\label{def:uniform-obligation}
A response obligation is \emph{uniform} when every $Q_p$ is the same
relation $Q$.
\end{definition}

The obligation is part of the declared spec.  Response Availability is only
as strong as $Q_p$: if a single output answers at every input, every spec
meets condition~(1) of Theorem~\ref{thm:complete-cap-response}.

A continuation of a partitioned history is \emph{within $p$'s side} when
the system produces it, it adds only history events of $p$'s side, and none of its
receives takes a message sent from another side.  A same-side continuation
for $p$ (Definition~\ref{def:local-continuation}) is one within $p$'s side
that adds no input event.

\begin{definition}[Response Availability]
\label{def:response-availability}
A system has \emph{Response Availability} for $(Q_p)_{p\in\mathcal N}$ when, in every
partitioned history it produces, each participant $p$, ending with current
input $i$, satisfies two conditions.
\begin{enumerate}
  \item \emph{Answering.}  In some continuation within $p$'s side that
        adds no input event at $p$, $p$ emits some output in $Q_p(i)$.
  \item \emph{Input acceptance.}  For every $j\InputGe i$, in some
        continuation within $p$'s side, $p$ records $\mathsf{input}(j)$.
\end{enumerate}
\end{definition}

Input acceptance says that $p$ never rules out an \indesc of its current
input.  It excludes specs whose answers require ruling out such an input, such
as the election under the equality order (Figure~\ref{fig:nonmonotone-frame}),
and it matches Gilbert and Lynch's requirement that nodes accept every
request.

\begin{definition}[Response Partition Consistency]
\label{def:response-partition-consistency}
A system has \emph{Response Partition Consistency} when, in every partitioned
history it produces and at every common \indesc $j$ of the participants'
current inputs, some $u\in\fiber{j}$ \outrel every output emitted in the
history.
\end{definition}

The converse of Theorem~\ref{thm:complete-cap} and
Corollary~\ref{cor:complete-cap-any-system} use only this part of Partition
Consistency, so both also hold with Response Partition Consistency in place of
Partition Consistency.  Confluence at $j$ implies a common \outnoun at $j$
(Corollary~\ref{cor:reconciliation-any-basis}).  The two coincide for a system
with basis $\mathcal B_P$ and no \commitment events, whose participant runs
are canonical runs, to which Lemma~\ref{lem:reconciliation} applies
(\S\ref{sec:distributed-consistency}).  With other \commitments they can
differ: participants may prune their states so that they cannot be
reconciled, while their emitted outputs still have a common \outnoun.

A system is \emph{replicated} (SPMD) when its obligation is uniform and every
participant runs the same strategy, which depends only on the participant's
initial input and local history, not on its identity or the number of
participants.

\begin{lemma}[Locality]
\label{lem:locality}
For each participant $p$ in some set of participants, let $\ell_p$ be a local
history that $p$'s
strategy can produce from initial input $i_p$ and that contains no receive.
Placing each $\ell_p$ on its own side, with at least two sides and
any further participants
idle, gives a partitioned history that the system produces.  If the system is
replicated, a receive-free local history that the strategy can produce for one participant
it can produce for any participant with the same initial input, in any
participant set.
\end{lemma}

\begin{proof}
Whether a history is produced is checked one history event at a time against each
participant's own local prefix, as fixed before
Definition~\ref{def:local-continuation},
and there are no receives to match.
Replication makes that check independent of the participant.
\end{proof}

\subsubsection{Response interfaces}

\begin{definition}[Response interface, self-continuing]
\label{def:response-interface}
For a uniform obligation $Q$, a \emph{response interface} is an explicit
relation $T\subseteq Q$ that is total on $I$ and \emph{self-continuing}: if
$(i,o)\in T$ and $i\InputOrd j$, then some $u\in T(j)$ \outrel $o$.
\end{definition}

Lifted Monotonicity for $T$ is that of Definition~\ref{def:lm-for-A}, and
the reconciling output may lie outside $T$.  With $T=P$ it is the
common-extension condition of Proposition~\ref{prop:common-expansion-form}, and
by Lemma~\ref{lem:reconciliation} it is confluence of the canonical runs whose
emitted rows lie in $T$.

A system can always serve a response interface.  Consider the canonical system
for $T$ (Section~\ref{sec:complete-cap}).  It emits, by the canonical output
event, any output of $T$ in the current fiber.  Each participant's emitted
outputs form a chain $o_1\OutputOrd\cdots\OutputOrd o_m$, so by
Lemma~\ref{lem:run-residual} its table is
$P[\mathrm{in}\InputGe i,\ \mathrm{out}\OutputGe o_m]$, or
$P[\mathrm{in}\InputGe i]$ before any emission.  Self-continuation applied to
the row of $o_m$, or totality of $T$, gives some $u\in T(i)$ in that table.
So every participant's fiber intersects $T$, and Availability for $T$ gives
Answering.  The same argument at an input $k\InputGe i$ leaves a row at $k$,
which gives Input acceptance.

\begin{theorem}[Complete CAP with chosen responses]
\label{thm:complete-cap-response}
For an \spectype $\Spec$ and a uniform obligation $Q$,
the following are equivalent:
\begin{enumerate}
  \item some replicated system has Response Availability for $Q$ and
        Response Partition Consistency;
  \item $\Spec$ satisfies Lifted Monotonicity for some response interface $T\subseteq Q$;
  \item for some response interface $T\subseteq Q$, some system has
        Availability for $T$ and Partition Consistency.
\end{enumerate}
\end{theorem}

\begin{proof}
(2$\Leftrightarrow$3) is Theorem~\ref{thm:complete-cap} and
Corollary~\ref{cor:complete-cap-any-system} with $E=T$.

(2$\Rightarrow$1.)  We take the canonical system for $T$.  It is replicated, and by the
paragraph above it has Answering and Input acceptance.  The outputs emitted in
a history are rows of $T$ at inputs that every common \indesc $j$ of the
current inputs \inrel, so Lifted Monotonicity for $T$ gives Response Partition
Consistency.

(1$\Rightarrow$2.)  We build $T$ from the outputs that participants emit
before receiving any message, which we call \emph{isolated} emissions.  Let
$T$ be the set of rows $(i,o)$ with $o\in Q(i)$ emitted in an isolated
emission in some partitioned history of the system; such a history and
participant are a \emph{witness} for $(i,o)$.  Then $T\subseteq Q$.  The
proof uses two facts about witnesses.  First, consider a participant alone on
its side.  Every continuation that Response Availability supplies adds only
its own history events, which in a partitioned history include no receive, so
the emission they yield is isolated.  Second, by Lemma~\ref{lem:locality}, a
witness's local history up to the emission can be replayed by a participant
alone on its side, with an idle participant forming a second side.

\emph{Totality:} for each $i$, Answering in the empty history, for a
participant with initial input $i$ alone on its side, gives a row of $T$ at
$i$.  \emph{Self-continuation:} let $(i,o)\in T$ and $j\InputGe i$.  We replay
the witness and record $\mathsf{input}(j)$ by Input acceptance, truncating
just after it.  Answering then gives an emission of some $u\in Q(j)$.  By
Lemmas~\ref{lem:canonical-equivalent} and~\ref{lem:run-residual}, the table
before that emission is contained in $P[\mathrm{in}\InputGe j,\
\mathrm{out}\OutputGe o]$.  So $o\OutputOrd u$ and $(j,u)\in T$.
\emph{Lifted Monotonicity for $T$:} let rows $(i_r,o_r)$ of $T$ have a common
\indesc $j$.  We replay each witness on its own participant, truncated after
the emission, and add a participant with initial input $j$.  By
Lemma~\ref{lem:locality} this is a partitioned history of the system.
Response Partition Consistency, applied to that history, gives
$u\in\fiber{j}$ \outrelpart every $o_r$.  If there are no rows, totality of
$P$ gives some $u\in\fiber{j}$.
\end{proof}

\begin{remark}[Without replication]
For a system that is not replicated, fix the participant set.  Let $T_p$ be
the rows $p$ emits in histories where it is alone on its side.  The same proof
gives one response interface $T_p\subseteq Q_p$ per participant.  Each
strategy sees only its own participant's initial input, and each participant
is replayed once.  So consistency is required only for families that draw on
each $T_p$ at most once.
\end{remark}

\subsubsection{Gilbert--Lynch with real-time precedence}

Section~\ref{sec:complete-cap} derives the Gilbert--Lynch theorem from
Corollary~\ref{cor:complete-cap-any-system}, with real-time precedence supplied
as input.  Here precedence comes from real time, and consistency is required
only at the precedence an execution realizes.

\begin{lemma}[Forced answer]
\label{lem:forced-answer}
Let $i\InputOrd j$, and suppose no output in $Q_q(i)$ has an \outnoun in
$\fiber{j}$.  In a system with Answering for $(Q_p)$
(Definition~\ref{def:response-availability}(1)), consider a partitioned
history that the system produces, in which $q$ is alone on its side,
$q$'s current input is $i$, and $j$ \inrel every participant's current input.
Then some continuation within $q$'s side that adds only history events of $q$,
and no input event, has emitted outputs with no common \outnoun in
$\fiber{j}$.
\end{lemma}

\begin{proof}
Answering gives such a continuation in which $q$ emits some $o\in Q_q(i)$.  No
current input changes, so $j$ still \inrel all of them, and no output of
$\fiber{j}$ \outrel $o$.
\end{proof}

\paragraph{The register.}
An input is a finite operation history: a set of read and write invocations at
replicas with a strict partial order of real-time precedence.  One history
\inrel another when it contains all of its operations and precedence edges.
An output is a response map assigning values to reads and acknowledgments to
writes, ordered by extension.  A linearization of $h$ is a total order of its
operations extending precedence~\cite{herlihy1990linearizability} in which
each answered read returns the last preceding write, or the initial value
$v_0$.  The fiber $\fiber{h}$ contains the maps explained by some
linearization of $h$.  The obligation $Q_p(h)$ consists of the maps in
$\fiber{h}$ that answer every read and acknowledge every write invoked at $p$.
Every finite history has a linearization, so $Q_p$ is total, and it is upward
closed within fibers.  Real-time precedence belongs to the input: the
environment invokes $b$ after seeing $a$'s response, so it supplies the edge
from $a$ to $b$ (\S\ref{sec:election-story}).

Real time is the order of history events in an execution.  An
\emph{interleaving} $\lambda$ of a distributed history is a linear order of
its history events consistent with happens-before.  No strategy observes it.
The \emph{invocation} of an operation is the first input event at its replica
whose input contains it.  Its \emph{response} is the first output event at
that replica whose map answers it.  For a prefix $C$ of $\lambda$ (a
consistent cut), the \emph{realized history} $h_C$ contains the operations
invoked in $C$, with an edge from $a$ to $b$ when $a$'s response precedes
$b$'s invocation in $\lambda$.  If $C'\subseteq C$ is also a prefix, $h_C$
\inrel $h_{C'}$.  The realized history need not be a common \indesc of the
participants' current inputs: an input event may introduce two operations with
an edge between them before either responds.

A \emph{client history} has two participants with empty initial inputs, and
each of its input events adds one operation, invoked at its participant after
all of that participant's earlier operations are answered.

\begin{corollary}[Gilbert--Lynch]
\label{cor:gilbert-lynch}
No system for the linearizable register has Answering for $(Q_p)$ together
with the following, both required only of client histories: for every
partitioned client history it produces, interleaving $\lambda$, and prefix $C$
of $\lambda$ such that $h_C$ is a common \indesc of the current inputs at $C$,
the outputs emitted in $C$ have a common \outnoun in $\fiber{h_C}$.
\end{corollary}

\begin{proof}
Let replicas $p$ and $q$ start with empty inputs on separate sides, and let
the environment invoke a write of $\dog\neq v_0$ at $p$.  Answering gives a
continuation in which $p$ acknowledges it.  The environment then invokes a read
$r_q$ at $q$, after the acknowledgment in $\lambda$.  The history is a client
history.  No message crosses the partition, so $q$'s current input is
$i_q=\{r_q\}$.  The realized history $j=h_C$ contains the write, $r_q$, and an
edge from the write to $r_q$, and it is a common \indesc of both current
inputs.  Every map in $Q_q(i_q)$ answers $r_q$ with $v_0$, and every
linearization of $j$ places the write first.  So no output in $Q_q(i_q)$ has
an \outnoun in $\fiber{j}$.  Lemma~\ref{lem:forced-answer} gives a
continuation adding no input event.  The realized history stays $j$, and the
emitted outputs have no common \outnoun in $\fiber{j}$.
\end{proof}

A system with Response Availability, and with a common \outnoun at $h_C$ for
every prefix $C$ of every partitioned history, satisfies these hypotheses, so
the corollary excludes it too.  The register's obligation differs by replica, so
Theorem~\ref{thm:complete-cap-response} does not apply as stated.  The
corollary yields Gilbert and Lynch's theorem in their asynchronous network
model~\cite{gilbert2002cap}.

\begin{proposition}[Gilbert--Lynch, asynchronous model]
\label{prop:gl-async}
In the asynchronous network model with at least two nodes, no algorithm for a
read/write register is both available (a request pending at a node is
answered in some extension of the execution that adds no request and loses
every message) and atomic (every execution with well-formed clients has a
linearization).
\end{proposition}

\begin{proof}[Proof sketch]
We translate the algorithm into a system with the canonical basis.  Participant
$p$ runs node $p$.  An invocation becomes the input event whose input is $p$'s
view: the operations invoked at $p$, with no precedence.  A response becomes an
output event emitting $p$'s response map, and $p$ may also re-emit its current
map.  With every message lost, $p$'s history is node $p$'s run in the
execution where it is alone.  Atomicity of that execution puts the map in the
fiber at the view, so the system can perform its output events under the
canonical basis.  Availability, applied to each unanswered operation of the
view, gives Answering.  An interleaving of a client history is an execution of
the algorithm whose real-time order is the precedence of $h_C$.  So atomicity
gives a common \outnoun in $\fiber{h_C}$, and the corollary excludes the
algorithm.
\end{proof}

Gilbert and Lynch require availability and atomicity in all fair executions,
including those in which messages are lost, in Lynch's asynchronous network
model~\cite[Theorem~1]{gilbert2002cap}\cite[Chapter~8]{lynch1996distributed}.
Both imply the forms in the proposition.  In their model any message may be
lost and client requests are inputs, so a finite execution extends to a fair
execution with no new request in which every message is lost, as in their
own proofs.  Availability answers the pending request at some finite point of
that extension, and the prefix up to that response is the extension the
proposition requires.  Atomicity of the fair extension gives atomicity of the
finite execution, since atomicity is a safety property and so holds of every
prefix of an atomic execution~\cite[Chapter~13]{lynch1996distributed}.
